\documentclass[11pt,a4paper]{article}
\usepackage[T1]{fontenc}
\usepackage[utf8]{inputenc}
\usepackage{amsmath,amssymb,xcolor}
\usepackage[margin=25mm]{geometry}
\usepackage[colorlinks=true,urlcolor=blue]{hyperref}
\definecolor{nexorablue}{HTML}{2459B5}
\setlength{\parindent}{0pt}
\setlength{\parskip}{7pt}
\setlength{\emergencystretch}{2em}
\renewcommand{\familydefault}{\sfdefault}
\begin{document}

{\small\bfseries\color{nexorablue} NEXORA CPGE}\par
{\small MPSI / NEX-ENS-02 / Corrigé}\par
{\LARGE\bfseries Décomposer une différence}\par
\vspace{0.25em}\hrule\vspace{0.6em}

\section*{Équivalence}
Être dans le membre de gauche signifie être dans $A$ et manquer au moins l’un des ensembles $B$ ou $C$. En distribuant la conjonction, on obtient l’appartenance au membre de droite.

\section*{Intersection des deux morceaux}
Leur intersection est $A\setminus(B\cup C)$. Pour $A=\{1\}$ et $B=C=\varnothing$, ce morceau n’est pas vide. La réunion est disjointe exactement quand $A\subset B\cup C$.


\end{document}
